\documentclass[10pt,a4paper]{amsart}
\usepackage[margin=19mm]{geometry}
\usepackage{amsmath,amssymb,mathtools}
\usepackage[scr=rsfs,cal=euler]{mathalfa}
\usepackage{xfrac}
\usepackage[unicode,hidelinks,bookmarks=false]{hyperref}
\newcommand{\ie}{i.e.\ }
\newcommand{\cf}{cf.\ }
\newcommand{\resp}{respectively\ }
\newcommand{\Subsection}{\subsection}
\providecommand{\nobreakdash}{\nobreak\hbox{-}}
\setlength{\emergencystretch}{2em}
\multlinegap0pt
\usepackage{bm}
\usepackage{amssymb}
\newtheorem{proposition}{Proposition}
\newcommand{\mtrx}[1]{\mathbf{#1}}
\newcommand{\pth}[1]  {\left( #1 \right)}
\title{An all-topology two-fluid model for two-phase flows derived through Hamilton's Stationary Action Principle}
\author{Ward Haegeman (CMAP), Giuseppe Orlando (CMAP), Samuel Kokh (MDLS), Marc Massot (CMAP)}
\date{Source: arXiv:2509.26298v2 (CC BY 4.0)}
\begin{document}
\maketitle

\noindent\emph{Source and licence.} An all-topology two-fluid model for two-phase flows derived through Hamilton's Stationary Action Principle. Version arXiv:2509.26298v2 (CC BY 4.0), distributed by arXiv under the Creative Commons Attribution 4.0 International licence, http://creativecommons.org/licenses/by/4.0/, as recorded in the arXiv licence metadata for this exact version and shown on the abstract page https://arxiv.org/abs/2509.26298v2. Full licence notice: \texttt{licenses/arxiv-2509.26298-CC-BY-4.0.txt}.\par\smallskip

\noindent\textit{Editorial scope.} Counted unit: the article's proposition prop:hyperbolicity (source0.tex lines 468--469). The article's complete original proof of it is delivered with it (source0.tex lines 470--511, one \texttt{proof} environment of 42 lines). No mathematics is written, repaired or re-derived: the statement and the proof are the article's own lines, in the article's own notation, and no retained line is substituted or re-flowed.  No further source-local context is needed: every cross-reference of every retained line resolves inside the two delivered spans.  Nothing was omitted from any retained span, so the package carries no omission record.  No retained line contains a citation, so the wrapper carries no bibliography and no bibliography entry.  No retained line contains a figure environment, an image-inclusion command or drawing code, so no image is delivered and the package declares no asset dependency.  Editorial formatting only, in addition: an AMS \texttt{amsart} preamble that restates the article's own declarations and macros used by the retained lines, namely the environments \texttt{proposition} on separate counters, together with the standard packages those lines need.  The retained environments print in this document as Proposition 1 in order of appearance, whereas the article prints the same items with the numbering of its own full document.  The complete native source of the article as distributed in the arXiv e-print for this exact version is bundled unchanged as \texttt{sources/186076/source0.tex}.
\begin{proposition}[Hyperbolicity]\label{prop:hyperbolicity} The system admits real eigenvalues which are $u$, $u_k$ and $u_k\pm c_k$, with $k=1,2$, and where $c_k=\partial p_k/\partial{\rho_k}_{|s_k}$, denotes the phasic sound velocity. The corresponding eigenvectors span $\mathbb{R}^7$, and the system is hyperbolic, under the non-resonance condition $u\neq u_k\pm c_k$, for $k=1,2$.
\end{proposition}

\begin{proof}
Through a change of variables, the system is expressed in terms of the variable $\bm{U}=(\alpha_1,u_1,p_1,s_1,u_2,p_2,s_2)^\top$, and writes in the following quasilinear form,
\begin{equation}\label{eq:quasilinear_form}
\partial_t\bm{U} + \mtrx{M}(\bm{U})\partial_x\bm{U}=\bm{0},
\end{equation}
with
\begin{equation}
\mtrx{M}(\bm{U}) = 
\begin{bmatrix}
u & \bm{0}_{1\times3} & \bm{0}_{1\times3} \\
\bm{z}_{\alpha,1}  & \mtrx{M}_1 & \bm{0}_{3\times3}\\
\bm{z}_{\alpha,2}  & \bm{0}_{3\times3} & \mtrx{M}_2
\end{bmatrix},
\end{equation}
and
\begin{equation}
\mtrx{M}_k = \begin{bmatrix}
u_k & 1/\rho_k & 0 \\
\rho_k c_k^2 & u_k & 0\\
0 & 0 & u_k
\end{bmatrix},
\qquad 
\bm{z}_{\alpha,k} = 
\begin{bmatrix}
\tfrac{p_1-p_2}{\rho} \\ (-1)^k\tfrac{\rho_k c_k^2}{\alpha_k}(u-u_k)\\ 0
\end{bmatrix}.
\end{equation}
The block-triangular structure of $\mtrx{M}(\bm{U})$ allows for an immediate computation of the eigenvalues. When the non-resonance condition is satisfied, the matrix $\mtrx{M}(\bm{U})$ is diagonalizable and a matrix of right eigenvectors $\mtrx{R}(\bm{U})$ is given by
\begin{equation}\label{eq:eigenvectors}
\mathbf{R}(\bm{U})= 
\begin{bmatrix}
1 & 0 & 0 & 0 & 0 & 0 & 0 \\
\tfrac{-Y_2 W}{(u_1-u)^2-c_1^2}\pth{\tfrac{p_1-p_2}{\rho} - \tfrac{c_1^2}{\alpha_1}} & 1 & 1 & 0 & 0 & 0 & 0 \\
\tfrac{\rho_1 c_1^2}{(u_1-u)^2-c_1^2}\pth{\tfrac{p_1-p_2}{\rho} - \tfrac{Y_2^2 W^2}{\alpha_1}} & \rho_1 c_1 & -\rho_1 c_1 & 0 & 0 & 0 & 0 \\
0 & 0 & 0 & 1 & 0 & 0 & 0\\
\tfrac{Y_1 W}{(u_2-u)^2-c_2^2}\pth{\tfrac{p_1-p_2}{\rho} + \tfrac{c_2^2}{\alpha_2}} & 0 & 0 & 0 & 1 & 1 & 0\\
\tfrac{\rho_2 c_2^2}{(u_2-u)^2-c_2^2}\pth{\tfrac{p_1-p_2}{\rho} + \tfrac{Y_1^2 W^2}{\alpha_2}}  & 0 & 0 & 0 & \rho_2 c_2 & -\rho_2 c_2 & 0\\
0 & 0 & 0 & 0 & 0 & 0 & 1
\end{bmatrix}.
\end{equation} 
The columns of $\mtrx{R}(\bm{U})$ correspond to eigenvectors $u$, $u_1+c_1$, $u_1-c_1$ ,$u_1$, $u_2+c_2$, $u_2-c_2$ and $u_2$  respectively.
\end{proof}

\end{document}
